% GENERATED FILE. Edit canonical lessons, then run npm run generate:books.
% canonical-source-hash: fnv1a64:9f7eb64fa27d8974
% canonical-lessons: RU-C177-obnimat, RU-C177-tselovat, RU-C177-sporit, RU-C177-ssoritsya, RU-C177-miritsya

\chapter{To hug, To kiss, To argue, To quarrel, To make up}
\label{ch:ru-to-hug-to-kiss-to-argue-to-quarrel-to-make-up}
\hlchaptermodality{177}

\begin{chapteropening}
\textbf{By the end of this chapter:} \emph{I can say to hug, to kiss, to argue, to quarrel and to make up.}

Complete the last lesson of chapter 177: I can say to hug, to kiss, to argue, to quarrel and to make up.
\end{chapteropening}

\section[\texorpdfstring{\ru{обнимать} (obnimát)}{обнимать (obnimát)}]{\ru{обнимать} (obnimát) — to hug}

\label{lesson:RU-C177-obnimat}



\begin{quote}
Before the new one: say the Russian for to rest, then the Russian for to get tired.
\end{quote}

\subsection*{You'll want to know: \ru{обнимать}}
\textbf{\ru{обнимать}} (\emph{obnimát}) — \textquotedblleft{}to hug\textquotedblright{}.

\begin{grammarlens}[title={What you've built}]
One more word for this chapter.
\end{grammarlens}

\subsection*{Your turn}
\begin{itemize}
\raggedright
  \item \emph{Say it:} \emph{obnimát}
  \item \emph{Say it:} \emph{obnimát}, once more
  \item \emph{Say it:} say \emph{obnimát}
  \item \emph{Recall:} say the Russian for to rest, then the Russian for to get tired, then say \emph{obnimát} again
\end{itemize}

\begin{tcolorbox}[breakable,colback=teal!4,colframe=teal!35!black,title={Before you move on}]
What does \ru{обнимать} mean? (\textquotedblleft{}To hug\textquotedblright{}.) Say it once more.
\end{tcolorbox}
\section[\texorpdfstring{\ru{целовать} (tselovát)}{целовать (tselovát)}]{\ru{целовать} (tselovát) — to kiss}

\label{lesson:RU-C177-tselovat}



\begin{quote}
Before the new one: say the Russian for to get tired, then the Russian for to hug.
\end{quote}

\subsection*{You'll want to know: \ru{целовать}}
\textbf{\ru{целовать}} (\emph{tselovát}) — \textquotedblleft{}to kiss\textquotedblright{}.

\begin{grammarlens}[title={What you've built}]
One more word for this chapter.
\end{grammarlens}

\subsection*{Your turn}
\begin{itemize}
\raggedright
  \item \emph{Say it:} \emph{tselovát}
  \item \emph{Say it:} \emph{tselovát}, once more
  \item \emph{Say it:} say \emph{tselovát}
  \item \emph{Recall:} say the Russian for to get tired, then the Russian for to hug, then say \emph{tselovát} again
\end{itemize}

\begin{tcolorbox}[breakable,colback=teal!4,colframe=teal!35!black,title={Before you move on}]
What does \ru{целовать} mean? (\textquotedblleft{}To kiss\textquotedblright{}.) Say it once more.
\end{tcolorbox}
\section[\texorpdfstring{\ru{спорить} (spórit)}{спорить (spórit)}]{\ru{спорить} (spórit) — to argue}

\label{lesson:RU-C177-sporit}



\begin{quote}
Before the new one: say the Russian for to hug, then the Russian for to kiss.
\end{quote}

\subsection*{You'll want to know: \ru{спорить}}
\textbf{\ru{спорить}} (\emph{spórit}) — \textquotedblleft{}to argue\textquotedblright{}.

\begin{grammarlens}[title={What you've built}]
One more word for this chapter.
\end{grammarlens}

\subsection*{Your turn}
\begin{itemize}
\raggedright
  \item \emph{Say it:} \emph{spórit}
  \item \emph{Say it:} \emph{spórit}, once more
  \item \emph{Say it:} say \emph{spórit}
  \item \emph{Recall:} say the Russian for to hug, then the Russian for to kiss, then say \emph{spórit} again
\end{itemize}

\begin{tcolorbox}[breakable,colback=teal!4,colframe=teal!35!black,title={Before you move on}]
What does \ru{спорить} mean? (\textquotedblleft{}To argue\textquotedblright{}.) Say it once more.
\end{tcolorbox}
\section[\texorpdfstring{\ru{ссориться} (ssóritsya)}{ссориться (ssóritsya)}]{\ru{ссориться} (ssóritsya) — to quarrel}

\label{lesson:RU-C177-ssoritsya}



\begin{quote}
Before the new one: say the Russian for to kiss, then the Russian for to argue.
\end{quote}

\subsection*{You'll want to know: \ru{ссориться}}
\textbf{\ru{ссориться}} (\emph{ssóritsya}) — \textquotedblleft{}to quarrel\textquotedblright{}.

\begin{grammarlens}[title={What you've built}]
One more word for this chapter.
\end{grammarlens}

\subsection*{Your turn}
\begin{itemize}
\raggedright
  \item \emph{Say it:} \emph{ssóritsya}
  \item \emph{Say it:} \emph{ssóritsya}, once more
  \item \emph{Say it:} say \emph{ssóritsya}
  \item \emph{Recall:} say the Russian for to kiss, then the Russian for to argue, then say \emph{ssóritsya} again
\end{itemize}

\begin{tcolorbox}[breakable,colback=teal!4,colframe=teal!35!black,title={Before you move on}]
What does \ru{ссориться} mean? (\textquotedblleft{}To quarrel\textquotedblright{}.) Say it once more.
\end{tcolorbox}
\section[\texorpdfstring{\ru{мириться} (mirítsya)}{мириться (mirítsya)}]{\ru{мириться} (mirítsya) — to make up (after a quarrel)}

\label{lesson:RU-C177-miritsya}



\begin{quote}
Before the new one: say the Russian for to argue, then the Russian for to quarrel.
\end{quote}

\subsection*{You'll want to know: \ru{мириться}}
\textbf{\ru{мириться}} (\emph{mirítsya}) — \textquotedblleft{}to make up (after a quarrel)\textquotedblright{}.

\begin{grammarlens}[title={What you've built}]
One more word for this chapter.
\end{grammarlens}

\subsection*{Your turn}
\begin{itemize}
\raggedright
  \item \emph{Say it:} \emph{mirítsya}
  \item \emph{Say it:} \emph{mirítsya}, once more
  \item \emph{Say it:} say \emph{mirítsya}
  \item \emph{Recall:} say the Russian for to argue, then the Russian for to quarrel, then say \emph{mirítsya} again
\end{itemize}

\begin{tcolorbox}[breakable,colback=teal!4,colframe=teal!35!black,title={Before you move on}]
What does \ru{мириться} mean? (\textquotedblleft{}To make up (after a quarrel)\textquotedblright{}.) Say it once more.
\end{tcolorbox}
